\section{Hyperparameters}

\Cref{tab:hyperparams} reports the main training hyperparameters shared by the
deep learning models.
The empirical design holds these settings fixed across LSTM, CS-Gated, and
TFT wherever possible so that performance differences can be interpreted as
differences in objective function and architecture rather than differences in
training budget.
The choices are intentionally conservative for a 324-week panel: the hidden
dimension is kept moderate, dropout and weight decay are non-trivial, and
early stopping is based on validation performance rather than test-period
feedback.

\begin{table}[t]
\centering
\caption{Training hyperparameters.}
\label{tab:hyperparams}
\small
\begin{tabular}{ll}
\toprule
\thead{Parameter} & \thead{Value} \\
\midrule
Lookback window $L$         & 8 weeks \\
Hidden dimension            & 64 \\
Attention heads             & 2 \\
LSTM layers                 & 1 \\
CS-Gated layers             & 3 \\
Dropout                     & 0.30 \\
Learning rate               & 0.001 (Adam) \\
Weight decay                & $10^{-4}$ \\
Batch size                  & 128 sequences \\
Grad.\ accumulation steps   & 4 \\
Effective batch (assets)    & $86 \times 4 = 344$ \\
ListNet temperature $\tau$  & 1.0 \\
Warmup epochs               & 10 \\
Epochs (max)                & 150 \\
Early stopping patience     & 20 \\
Number of deep-model seeds  & 32 \\
Number of tree-model seeds  & 32 \\
\midrule
Train / Valid / Test split  & 70\% / 15\% / 15\% \\
Test weeks                  & 49 \\
\bottomrule
\end{tabular}
\end{table}

\section{Information Sets}

We evaluate six information sets of increasing richness:

\begin{enumerate}
  \item \textbf{Price+Technical} (11 features): momentum and technical indicators only.
  \item \textbf{+Onchain} (16): adds on-chain metrics.
  \item \textbf{All} (33): all baseline market, macro, ETF, and prediction-market features.
  \item \textbf{+Trump} (38): \texttt{All} plus 5 Trump social media features.
  \item \textbf{+ETF} (37): \texttt{All} plus 4 granular Farside ETF features.
  \item \textbf{+Trump+ETF} (42): \texttt{All} plus both Trump and Farside ETF features.
\end{enumerate}

This design is therefore not intended as a simple ``more features is better''
progression.
Instead, it is structured to identify which types of information are most
usefully handled by which types of models.
The \texttt{+Onchain} configuration tests whether contemporaneous
fundamentals already contain enough information for cross-sectional ranking,
whereas \texttt{+Trump} and \texttt{+ETF} isolate two forms of
alternative data with more explicit temporal structure.

\section{Ensemble Strategy}

For each (model, feature-configuration) pair, we train 32 independent deep
learning runs with different random seeds and average their predicted rank
scores before portfolio construction.
The purpose of ensembling here is not only to improve headline performance,
but also to stabilize a training problem in which single-seed results can
vary substantially.
The averaging is performed at the score level, before assets are sorted into
deciles.
No seed is selected or discarded using test-set returns; each of the 32
trained networks contributes equally to the final ranking score.
This matters because the investment object is an ordering, not a calibrated
return forecast: averaging scores preserves the ordinal information learned
by different initializations while reducing idiosyncratic seed-level noise.

Traditional tree models are also averaged across 32 random states, while
linear models are deterministic.
This asymmetry reflects typical deployment practice for each model family,
but it also means that model comparisons should be interpreted along two
dimensions: best achievable ensemble performance, and the stability of an
individual training run.
We return to this distinction in \Cref{sec:analysis:ensemble}.
